\section{Channels, instruments, and the meaning of coherent access}
\label{sec:channels-foundations}
The passage from vectors to density operators is not merely a way to add
noise to an ideal circuit. It specifies which transformations remain
physical when a register is entangled with an unobserved reference.
Throughout this section all Hilbert spaces are finite-dimensional.

\subsection{Complete positivity and an environment representation}
A quantum channel $\Phi:\mathcal L(\mathcal H_A)\to
\mathcal L(\mathcal H_B)$ is linear, trace preserving, and completely
positive: $\Phi\otimes\operatorname{id}_R$ preserves positivity for
every reference space $\mathcal H_R$. Positivity on isolated inputs is
weaker. For example, transposition preserves eigenvalues, but partial
transposition of $\ket{\Omega}\bra{\Omega}$, with
$\ket{\Omega}=(\ket{00}+\ket{11})/\sqrt2$, has eigenvalues
$1/2,1/2,1/2,-1/2$. It cannot describe a physical local channel.

\begin{theorem}[Finite-dimensional channel representation]\label{thm:kraus-foundations}
A linear map is completely positive if and only if
\begin{equation}\label{eq:kraus-foundations}
 \Phi(A)=\sum_j K_j A K_j^\dagger.
\end{equation}
It is trace preserving exactly when $\sum_jK_j^\dagger K_j=I$.
In that case $V\ket{\psi}=\sum_jK_j\ket{\psi}\otimes\ket{j}$ is an
isometry and $\Phi(A)=\Tr_E(VAV^\dagger)$.
\end{theorem}
\begin{proof}
The forward representation is the finite-dimensional Choi theorem
\citep{choi1975}; its construction is short. Choose an orthonormal input
basis and form
$J_\Phi=\sum_{i,k}\ket{i}\bra{k}\otimes\Phi(\ket{i}\bra{k})$.
Complete positivity implies $J_\Phi\succeq0$, since it is the image
of the unnormalized maximally entangled projector under
$\operatorname{id}\otimes\Phi$. Write
$J_\Phi=\sum_j\ket{v_j}\bra{v_j}$ and define $K_j$ by
$\ket{v_j}=\sum_i\ket{i}\otimes K_j\ket{i}$.
Comparison of the $(i,k)$ blocks gives~\eqref{eq:kraus-foundations}.
Conversely, each term preserves positivity after tensoring with any
identity. Cyclicity of trace gives the trace-preservation condition,
which also implies $V^\dagger V=I$. Partial trace gives the last formula.
\end{proof}

The environment description explains decoherence without invoking an
unobserved classical measurement. For an isometry satisfying
$V\ket{0}=\ket{0}\ket{e_0}$ and $V\ket{1}=\ket{1}\ket{e_1}$,
discarding the environment multiplies the input $01$ matrix entry by
$\braket{e_1}{e_0}$. Orthogonal environment states erase that coherence;
identical environment states preserve it. The qubit dephasing channel
with visibility $v\in[0,1]$ has the explicit form
\begin{equation}\label{eq:dephasing-foundations}
 \mathcal D_v(\rho)=\frac{1+v}{2}\rho+
                   \frac{1-v}{2}Z\rho Z.
\end{equation}
It maps a Bloch vector $(x,y,z)$ to $(vx,vy,z)$, and
$\mathcal D_v\mathcal D_w=\mathcal D_{vw}$. Taking $v=e^{-\gamma}$
therefore gives visibility $e^{-\gamma t}$ after $t$ commuting phase
queries, exactly as assumed in Part~\ref{part:operational}.

\subsection{Outcome probabilities and post-measurement states}
A POVM $\{E_y\}$ specifies probabilities, but not the state after an
outcome. An instrument consists of completely positive maps
$\mathcal M_y$ whose sum is trace preserving. Its effect is
$E_y=\mathcal M_y^*(I)$; thus
\[
 p(y)=\Tr\mathcal M_y(\rho)=\Tr(E_y\rho),\qquad
 \rho_y=\mathcal M_y(\rho)/p(y)
\]
when $p(y)>0$. Distinct instruments can have the same POVM and different
state disturbances. In a restarted experiment the quantum output is
discarded, so only the POVM probabilities remain relevant. In a protocol
with quantum memory, that simplification is unavailable.

Any POVM has a direct isometric implementation:
$W\ket{\psi}=\sum_y\sqrt{E_y}\ket{\psi}\otimes\ket{y}$, for which
$W^\dagger W=I$. Measuring the last register realizes its effects.
This establishes physical realizability in principle. It does not make
the synthesis of $W$ free; a cost model must decide how to price it.

The dual channel satisfies
$\Tr[B\Phi(A)]=\Tr[\Phi^*(B)A]$ and is positive and unital.
Thus $0\preceq E\preceq I$ implies
$0\preceq\Phi^*(E)\preceq I$. This simple observation is the basis
for the distinguishability contraction proved below.

\subsection{Coherent access is more than access to samples}
Three interfaces must be distinguished: samples from a distribution;
preparation of a specified quantum state; and reversible access to a
preparation unitary, its inverse, and required controlled versions.
Each can supply different algorithmic information. A state-preparation
channel that discards its environment need not supply a coherent inverse.
An unspecified physical unitary channel also has no observable global
phase, while a controlled unitary turns that phase into a relative phase.

\begin{lemma}[The phase ambiguity of controlled access]
An interface giving only the channel $\rho\mapsto U\rho U^\dagger$
cannot specify a transformation that outputs the controlled unitary
$\ket{0}\bra{0}\otimes I+\ket{1}\bra{1}\otimes U$ for every chosen matrix
representative $U$.
\end{lemma}
\begin{proof}
$U=I$ and $U=-I$ define the same input channel. Their requested
controlled channels act as $I$ and $Z$ on the control register and
are distinguishable on $\ket{+}$. A transformation of the supplied
channel cannot distinguish two identical inputs.
\end{proof}

This elementary obstruction addresses a phase-specific target. Stronger
no-control results for an unknown unitary are established by
\citet{araujo2014}. A known circuit decomposition or a specified physical
reference can provide the missing structure. The controlled oracle in
the worked decision experiment is therefore a stated interface
assumption, not a consequence of ordinary sampling access.

\section{Distinguishability, information, and the choice of loss}
\label{sec:statistical-foundations}
Quantum states do not have one universally optimal measurement. An
optimal readout depends on the loss function, the prior or worst-case
criterion, the number of observations, and whether later experiments may
depend on earlier outcomes.

\subsection{The operational meaning of trace distance}
\begin{theorem}[Binary minimum-error discrimination]\label{thm:helstrom-foundations}
For states $\rho,\sigma$ with priors $\pi,1-\pi$, the minimum average
error over all measurements and binary decisions is
\begin{equation}\label{eq:helstrom-foundations}
 e_\pi^*(\rho,\sigma)=
 \frac{1-\|\pi\rho-(1-\pi)\sigma\|_1}{2}.
\end{equation}
For equal priors this is $(1-D_{\rm tr}(\rho,\sigma))/2$.
\end{theorem}
\begin{proof}
Merge all outcomes deciding $\rho$ into an effect $E$, with
$0\preceq E\preceq I$. Its success probability is
$(1-\pi)+\Tr(E\Delta)$, where
$\Delta=\pi\rho-(1-\pi)\sigma$.
The spectral decomposition maximizes the trace by taking $E$ to be
the projector onto the positive eigenspace of $\Delta$.
Since $\Tr\Delta=2\pi-1$,
$\Tr\Delta_+=(\|\Delta\|_1+2\pi-1)/2$. Subtraction from one
gives~\eqref{eq:helstrom-foundations}, the binary quantum detection
formula \citep{helstrom1969}.
\end{proof}

For traceless Hermitian $A$, the same spectral argument gives
$\max_{0\preceq E\preceq I}|\Tr(EA)|=\|A\|_1/2$.
Consequently trace distance is the largest possible difference in
measurement-event probabilities. It is also the maximum total
variation distance over POVMs: group the outcomes with positive
probability difference into one effect.

\begin{corollary}[Processing cannot improve single-copy discrimination]
For every channel $\Phi$,
\begin{equation}
 D_{\rm tr}(\Phi(\rho),\Phi(\sigma))
 \leq D_{\rm tr}(\rho,\sigma).
\end{equation}
\end{corollary}
\begin{proof}
Pull any output effect back through the positive unital map $\Phi^*$.
It becomes an admissible input effect, so the supremum of its probability
difference cannot exceed the input supremum.
\end{proof}

\subsection{Pure-state overlap and deterministic sample budgets}
For two pure states choose phases so that their overlap is
$q\in[0,1]$. In their two-dimensional span, write
$\ket{\phi}=q\ket{\psi}+\sqrt{1-q^2}\ket{\psi^\perp}$.
The difference of their rank-one projectors has eigenvalues
$\pm\sqrt{1-q^2}$. This proves~\eqref{eq:tracedistance} and gives,
for $n$ independent copies and an arbitrary collective measurement,
\begin{equation}\label{eq:pure-fixed-budget}
 e_n^*=\frac{1-\sqrt{1-q^{2n}}}{2}.
\end{equation}
For $0<q<1$ and $0<\delta<1/2$, the condition $e_n^*\leq\delta$
is equivalent to
\[
 n\geq\frac{\log[1/(4\delta(1-\delta))]}{-2\log q}.
\]
Every finite deterministic copy budget has positive error. This does not
contradict a sequential zero-error protocol with finite mean cost:
its stopping time can have an unbounded tail. In the symmetric
unambiguous construction of Proposition~\ref{prop:usd},
$\Pr(N>n)=q^n$ and $\E N=(1-q)^{-1}$.
The distinction between expected and deterministic resources is therefore
mathematical, not terminological.

For a POVM on a pure-state pair, the classical affinity
$\sum_y\sqrt{p_yq_y}$ is at least the overlap $q$.
Indeed,
\[
 q=\left|\sum_y\bra{\psi} E_y\ket{\phi}\right|
 \leq\sum_y\sqrt{\bra{\psi} E_y\ket{\psi}\,
                        \bra{\phi} E_y\ket{\phi}}.
\]
Unlike relative entropy, this quantity remains finite at the pure-state
boundary. Its decrement, rather than a divergent KL rate, controls the
noiseless cost converse in Theorem~\ref{thm:noiseless-converse}.

\subsection{Relative entropy measures accumulated evidence}
For finite probability laws $p,q$ define
$D_{\rm KL}(p\|q)=\sum_y p_y\log(p_y/q_y)$, with the usual support
convention. It is the expected log-likelihood increment under $p$.
The opposite direction generally has a different value. Under a common
classical channel $K$, the log-sum inequality applied separately at each
output gives
\[
 D_{\rm KL}(pK\|qK)\leq D_{\rm KL}(p\|q).
\]
For product observations, independence makes their log likelihoods
additive. Sequential designs preserve this conditional additivity even
when the observation rule is selected from the preceding transcript.

The quantum relative entropy
$D(\rho\|\sigma)=\Tr\rho(\log\rho-\log\sigma)$ is finite only
when the support of $\rho$ is contained in that of $\sigma$.
The measured relative entropy is
\[
 D_M(\rho\|\sigma)=\sup_{\{E_y\}}
 D_{\rm KL}\bigl(\{\Tr\rho E_y\}\|\{\Tr\sigma E_y\}\bigr).
\]
This is the evidence accessible from one measured copy, whereas the
operator expression need not be attainable by a single-copy readout.
Part~\ref{part:operational} makes that distinction quantitative in a
two-state family. Trace distance optimizes immediate binary decision
error; relative entropy optimizes the long-run directed evidence rate.
The optimizers need not coincide.

\subsection{Accessible classical information is limited by dimension}
Let a classical label $X=x$ select state $\rho_x$ with probability $p_x$,
and put $\bar\rho=\sum_xp_x\rho_x$. For a measurement outcome $Y$,
the Holevo bound is
\begin{equation}\label{eq:holevo-foundations}
 I(X:Y)\leq\chi:=S(\bar\rho)-\sum_xp_xS(\rho_x)
 \leq\log d,\qquad S(\rho)=-\Tr\rho\log\rho,
\end{equation}
where $d$ is the transmitted system dimension \citep{holevo1973}.
One way to see the first inequality is to write
$I(X:Y)=\sum_xp_xD_{\rm KL}(p_{Y|x}\|p_Y)$, apply measurement
data processing to each term, and use
$\sum_xp_xD(\rho_x\|\bar\rho)=\chi$.
The last inequality follows from $0\leq S(\rho_x)$ and
$S(\bar\rho)\leq\log d$.
Thus an $n$-qubit system, without preshared entanglement or additional
side information, conveys at most $n$ bits of accessible classical
information per transmitted system. This limits readout; it does not
rule out computing a useful small output from a large structured input.

\section{Sequential experiments with an explicit resource price}
\label{sec:general-sequential}
The information-per-cost principle does not depend on a quantum circuit.
It becomes a quantum resource statement only after the action set and
observation laws have been derived from the permitted physical operations.

\subsection{A controlled experiment and its stopped likelihood}
Let $h\in\{+,-\}$. At step $j$ a protocol chooses action $a_j$ from
the classical history, pays $c(a_j)\geq c_{\min}>0$, and observes an
outcome with conditional law $p_{h,a_j}$. It may randomize using
hypothesis-independent coins and stop at $N$, with finite expected
cost under both hypotheses. Assume common support and a uniform bound
on $|\log(p_{+,a}(y)/p_{-,a}(y))|$.
Its log likelihood is
\begin{equation}\label{eq:general-loglikelihood}
 Z_N=\sum_{j=1}^N\log\frac{p_{+,a_j}(Y_j)}{p_{-,a_j}(Y_j)}.
\end{equation}
The action-selection and stopping kernels cancel in the likelihood
ratio because they are the same functions of the observed history
under either hypothesis. They can change what is observed but cannot
create evidence without observations.

Define the directed rate ceilings
\[
 J_+=\sup_a\frac{D_{\rm KL}(p_{+,a}\|p_{-,a})}{c(a)},\qquad
 J_-=\sup_a\frac{D_{\rm KL}(p_{-,a}\|p_{+,a})}{c(a)}.
\]
\begin{theorem}[A costed sequential testing converse]\label{thm:general-converse}
If both decision errors are at most $0<\delta<1/2$, then
\begin{equation}\label{eq:general-cost-converse}
 \max_h\E_h\sum_{j=1}^N c(a_j)
 \geq\max\left\{\frac{k_\delta}{J_+},
                \frac{k_\delta}{J_-}\right\},\qquad
 k_\delta=(1-2\delta)\log\frac{1-\delta}{\delta}.
\end{equation}
The convention $k_\delta/0=+\infty$ expresses impossibility when one
directed rate is zero.
\end{theorem}
\begin{proof}
Finite expected cost gives $\E_hN<\infty$. Bounded increments justify
the stopped likelihood identity and interchange of summation and
expectation. Under $h=+$,
\[
 D_{\rm KL}(P_+^{\rm stop}\|P_-^{\rm stop})
 =\E_+Z_N
 =\E_+\sum_{j=1}^N
      D_{\rm KL}(p_{+,a_j}\|p_{-,a_j})
 \leq J_+\E_+\sum_{j=1}^Nc(a_j).
\]
Coarsen the terminal transcript to the decision. If its errors are
$\alpha,\beta\leq\delta$, data processing gives the lower bound
$\operatorname{kl}(1-\alpha,\beta)\geq
\operatorname{kl}(1-\delta,\delta)=k_\delta$.
The minimization over this rectangle follows by differentiating binary
KL in its two arguments while $1-\alpha>\beta$.
Reverse the hypotheses to obtain the $J_-$ inequality.
\end{proof}

This is a costed version of the information argument underlying
sequential likelihood-ratio testing and controlled sensing
\citep{wald1945,chernoff1959,nitinawarat2015}. It is a converse, not an
automatic achievability theorem. If different actions maximize the two
directions, an adaptive controller must select the appropriate action
using incomplete information about the true hypothesis. Reflection
symmetry makes that issue tractable in Theorem~\ref{thm:measured}.

\subsection{Finite confidence, asymptotic rates, and resource frontiers}
An asymptotic statement $C_\delta/\log(1/\delta)\to1/J$ describes
the leading coefficient for fixed experiment parameters. It does not
remove overshoot, integer-depth choices, or the learning cost of selecting
an action at finite confidence. In particular, two methods may have
ordered limiting rates and still require a separate finite-confidence
comparison. The cost interval $[L_\delta,U_\delta]$ records both a
converse and an implementable protocol.

For methods $A,B$, the condition $U_A<L_B$ certifies that $A$ wins
throughout the corresponding minimax cost intervals. The condition
$L_A>U_B$ certifies that it loses. Overlapping intervals are unresolved
by those bounds; they are not evidence of equal cost. Multiplying only
one oracle price can change the ordering without changing either
algorithm's query complexity.

Suppose an action consumes a resource vector $v(a)\in\R_+^m$, including
preparation, queries, measurement, and classical processing. A positive
price vector $w$ yields $c_w(a)=w^\top v(a)$ and rate
$J(w)=\sup_a d(a)/(w^\top v(a))$. Changing $w$ explores different
resource tradeoffs; there is no implementation-independent conversion
from queries to elapsed time. A protocol dominated in every expected
resource coordinate cannot become optimal for a strictly positive price
vector. The affine price $c_0+c_1t$ below is a transparent two-resource
specialization of this general viewpoint.

The order of limits is equally substantive. Fixed positive noise and
$\delta\downarrow 0$ retain bounded single-block likelihood increments
and produce logarithmic confidence cost. At zero noise, disjoint support
of selected outcomes can make those increments infinite and permit
finite-mean exact discrimination. A uniform result must address both
mechanisms rather than substitute zero noise into a full-rank formula.
